\input _template

\setlanguage{CS}

\Title{Číslování}

\MathcompsLink{cifrovani}

\Author{Patrik Bak}

\sec Úvod

Ukážeme si základní tipy a~triky používané v~úlohách s~číslicemi.

\sec Techniky

\begitems \style i
\i Číslo $\overline{a_1a_2\dots a_n}$ rozepíšeme pomocí dekadického zápisu, např. $\overline{ab}=10a+b$ nebo $\overline{abc}=100a+10b+c$.
\i Číslo se sekvencí číslic $x$ následovanou sekvencí číslic $y$ délky $k$ zapíšeme jako $x \cdot 10^k + y$, např. 4267 je pro $k=2$ rovno $42 \cdot 10^2 + 67$.
\i U~úloh s~číslicemi se vyplatí dívat se na zbytky po dělení~3 nebo~9, protože číslo dává stejný zbytek po dělení~3 nebo~9 právě když je tomu tak pro ciferný součet.
\i Další dobrá kritéria pro číslice jsou pro dělitelnost $2^n$ (posledních $n$ číslic tvoří číslo dělitelné~$2^n$).
\i A ještě jedno kritérium: číslo je dělitelné~11 právě když součet číslic na sudých místech dává stejný zbytek po dělení~11 jako součet číslic na lichých místech.
\enditems

\sec Úlohy

\Problem{0}{}{
	Najděte všechna čtyřciferná čísla $\overline{abcd}$ s~ciferným součtem~12 taková, že $\overline{ab}-\overline{cd}=1$.
}{
	Platí $\overline{ab}=\overline{cd}+1$. Rozeberte případy podle toho, zda při sčítání napravo dochází k~přechodu přes~10.
}{
}

\Problem{0}{}{
	Určete všechna přirozená čísla $n$, pro která platí
	$$
	n+p(n)=70,
	$$
	přičemž $p(n)$ označuje součin všech číslic čísla $n$.
}{
	Zapište $n=\overline{ab}$ a~rovnici přepište.
}{
}

\Problem{0}{}{
	Dvojciferné číslo $\overline{ab}$ nazveme \textit{nafukovatelné}, pokud z něj po přičtení 990-násobku vhodného jednociferného čísla získáme čtyřciferné číslo tvaru $\overline{axxb}$ s~nenulovou číslicí~$x$. Kolik nafukovatelných čísel existuje?
}{
	Zapište definici nafukovatelnosti pomocí rovnice $\overline{ab}+990y=\overline{axxb}$, kterou dále upravte.
}{
	Měli bychom dojít k~rovnici $9(y-a)=x$. Jenže~$x$ je nenulová číslice, to nám mnoho možností nedá.
}{
}

\Problem{0}{}{
	Zjistěte, jaké nejmenší kladné celé číslo lze vložit mezi dvojčíslí $20$ a $16$ tak, aby výsledné číslo bylo násobkem čísla $2016$.
}{
	Podívejme se na modulo~$9$.
}{
	Jelikož~$2016$ je dělitelné~9, tak ciferný součet vkládaného čísla musí být dělitelný~9, protože ciferný součet $20$ a~$16$ je~$9$.
}{
	Postupně zkoušejte vkládat $9, 18, \dots$.
}

\Problem{0}{}{
	Najděte všechna trojciferná čísla, která jsou součtem faktoriálů svých číslic.
}{
	Takové číslo nemůže mít příliš velké číslice. Jakou největší může mít?
}{
	Rozmyslíme, že největší číslice v~čísle může být 5. Musí tam vůbec nějaká být?
}{
	Rozmyslíme, že číslice~5 tam nutně musí být. Už to není mnoho možností, obzvlášť když se zamyslíme nad tím, co může být první číslice.
}{
}

\Problem{0}{}{
	Najděte všechna čtyřciferná čísla taková, že jsou druhou mocninou celého čísla, a~zároveň jejich první číslice je stejná jako druhá a~třetí je stejná jako čtvrtá.
}{
	Zapišme si naše číslo jako $\overline{aabb}$ a~rozepišme, čemu je rovno.
}{
	Vychází nám $\overline{aabb}=11(100a+b)$. Aby bylo toto čtvercem, potřebujeme, aby $100a+b$ bylo dělitelné~11.
}{
	Platí $100a+b=99a+(a+b)$, takže nutně $11 \mid a+b$. To už není mnoho možností.
}{
}

\Problem{0}{}{
	Najděte největší prvočíslo složené z~10 různých číslic.
}{
	Podívejte se na ciferný součet našeho čísla.
}{
}

\Problem{0}{}{
	Zdůvodněte, že každý palindrom se sudým počtem číslic je dělitelný~11.
}{
	Kritérium dělitelnosti~11.
}{
}

\Problem{0}{}{
	Číslo $123456789101112\dots 100$ postupně nahrazujeme ciferným součtem, dokud nezískáme jednociferné číslo. Jaké číslo nám zbyde?
}{
	Stačí najít ciferný součet a použít dělitelnost~9.
}{
	Pro spočítání ciferného součtu našeho čísla se zamysleme, kolikrát se každá z~číslic vyskytuje v~našem čísle.
}{
}

\Problem{0}{}{
	Existuje číslo dělitelné~11 složené z~číslic 1, 2, 3, 4, 5 a~6 tak, že každou použijeme právě jednou?
}{
	Použijeme kritérium dělitelnosti~11 -- chceme, aby rozdíl součtu číslic na sudých místech~$s_s$ a~na lichých místech~$s_l$ byl dělitelný~11. Zkuste najít možnosti pro tyto součty.
}{
	Uvědomíme si, že $s_s+s_l$ je ciferný součet našeho čísla, takže $1+2+3+4+5+6=21$. Z~toho najdeme možnosti pro dvojice $s_s$ a~$s_l$ a~ty zanalyzujeme.
}{
}

\Problem{0}{}{
	Mysli si trojciferné číslo, které neobsahuje nulu a má první a poslední číslici různou. Napiš ho pozpátku a odečti menší číslo od většího. Vzniklé číslo případně doplň zleva nulou na trojciferné, znovu přepiš pozpátku a tato dvě čísla sečti. Jaké nejmenší a~jaké největší číslo jsme mohli dostat?
}{
	Zapišme si číslo jako $\overline{abc}$. Co nám zbyde po první operaci?
}{
	V~dalším kroku máme $|\overline{abc}-\overline{cba}|=99|a-c|$. Vypište si tato možná čísla.
}{
}

\Problem{0}{}{
	Najděte všechna čísla s~první číslicí~6 taková, že po jejím odstranění vznikne 25-krát menší číslo.
}{
	Zapište si číslo jako $6 \cdot 10^k + x$, kde $k$ je počet číslic~$x$.
}{
	Sestavte si rovnici ze zadání: $6 \cdot 10^k + x = 25x$. Z~toho vyjádřete $x$ a~zjistěte, pro jaké $k$ bude $x$ celé číslo.
}{
}

\Problem{0}{}{
	Dokažte, že číslo $\underbrace{44\dots4}_{n+1}\underbrace{88\dots8}_{n}9$ je pro každé celé nezáporné číslo~$n$ čtvercem přirozeného čísla.
}{
	Rozepišme si číslo pomocí dekadického zápisu.
}{
	Připomeňte si, že číslo zapsané jako $k$ stejných číslic $c$ se dá vyjádřit jako $c \cdot \frac{10^k - 1}{9}$. Nezapomeňte každou část čísla vynásobit příslušnou mocninou desítky podle toho, kolik číslic za ní následuje.
}{
}

\Problem{0}{}{
	Je dáno přirozené číslo $n$. Určete poslední číslici čísla $n^2$, pokud víte, že~7 je jeho předposlední číslice.
}{
	Zapište číslo jako $10x+y$, kde $y$ je poslední číslice. Jak vypadá $n^2$?
}{
	Platí $n^2=100x^2+20xy+y^2$. Kdy může být 7 předposlední číslice?
}{
	Analyzujte vliv jednotlivých členů součtu $100x^2+20xy+y^2$ na místo desítek.
}{
	Hlavní idea je, že $100x^2$ nepřispívá na místo desítek vůbec a~$20xy$ přispívá sudou číslicí, takže $y^2$ musí mít lichou číslici na místě desítek.
}{
}

\Problem{1}{}{
	Existují dvě různé mocniny dvojky se stejným počtem číslic, z nichž by se jedna dala získat pouze přeuspořádáním číslic té druhé.
}{
	Podívejme se na modulo~$9$.
}{
	Zkoumejme mocniny dvojky modulo~9. Jednotlivé mocniny $2^1, 2^2, \dots$ jsou $2,4,8,7,5,1$ a~tak~dále. Za jakých podmínek $2^a$ a~$2^b$ dávají stejný zbytek po dělení~9?
}{
}

\Problem{1}{}{
	Pro která přirozená $n$ existuje $n$-ciferné číslo dělitelné $5^n$, které má všechny číslice liché?
}{
	Zkusme taková čísla konstruovat induktivně.
}{
    Vezměme si vyhovující číslo, které má $n$ číslic. Dokažte, že připojením jedné z~číslic $\{1,3,5,7,9\}$ na začátek čísla dostaneme vyhovující číslo s $n+1$ číslicemi.
}{
}

\Problem{1}{}{
	Uvažujme následující proces: Začneme s~libovolným přirozeným číslem $a_1$ a~vygenerujeme posloupnost $a_1,a_2,\dots$ tak, že $a_{n+1}$ dostaneme z~$a_n$ přilepením číslice různé od~9. Dokažte, že v~této posloupnosti bude nevyhnutelně existovat nekonečně mnoho složených čísel.
}{
	Postupujeme sporem. Co kdyby tam bylo jen konečně mnoho složených čísel? Rozmyslete si, jaké číslice vůbec dokážeme připojovat.
}{
	Určitě neumíme připojovat sudé číslice a~5. Zbývá 1, 3, 7. Zkuste vyloučit některé z~nich.
}{
	Idea je vyloučit 1 a~7. Těchto můžeme mít jen konečně mnoho kvůli tomu, že bychom porušili dělitelnost~3. Zbývají trojky. Můžeme po čase mít pouze ty?
}{
	Trik na trojky je rozepsat si číslo jako
	$$
	a\cdot 10^n + \underbrace{\overline{33\dots3}}_n=
	a\cdot 10^n + \frac{10^n-1}{3}.
	$$
	Druhou část umíme také napsat explicitně. Rozmyslete si, že by stačilo, aby $10^n-1$ bylo dělitelné~$3a$.
}{
}

\Problem{2}{}{
	Dokažte, že existuje nekonečně mnoho kladných celých čísel $n$ takových, že $n^2$ zapsané ve~čtyřkové soustavě obsahuje pouze číslice $1$ a~$2$.
}{
	Zkusme taková čísla konstruovat induktivně.
}{
	Trikem je, že ze tří kopií čísla umíme zkonstruovat nové, které vyhovuje. Nebojte se zesílit indukční předpoklad.
}{
}

\DisplayExerciseSolutions

\DisplayHints

\DisplayProblemSolutions

\bye
